\subsubsection{TAE Results}
\label{sec:tae_results}
We present here results and ablations from important design decisions for the temporal autoencoder (TAE).
For evaluation, we report the reconstructed peak signal-to-noise ratio (PSNR), structural similarity (SSIM)~\citep{wang2004image}, and Fr\'echet Inception distance (FID)~\citep{heusel2017gans} of video clips split from the training set, with 2s, 4s, 6s, and 8s duration, each with 200 examples.
We also measure the same metrics on a validation split of the image training set.
For video reconstruction evaluation, metrics are averaged over video frames.

\noindent\textbf{Qualitative Results.}
We show sample reconstructions from our \taeShort in~\cref{fig:tae_examples} with frames from the original video and the reconstruction after the \taeShort encoder and decoder.
We observe that the \taeShort can reconstruct the video frames while preserving visual detail.
The \taeShort reconstruction quality decreases for high frequency spatial details in images and video frames, and fast motion in videos.
When both high frequency spatial details and large motion are present in a video, this can lead to a loss in detail, as can be seen in the examples in~\cref{fig:tae_examples}, where fine details are smoothed out in the reconstruction.

\begin{figure}[t!]
  \centering
  \begin{subfigure}[b]{0.49\textwidth}
      \centering
      \includegraphics[width=1\linewidth]{figures/foundation_model/cat_sbs.png}
  \end{subfigure}
  \begin{subfigure}[b]{0.49\textwidth}
      \centering
      \includegraphics[width=1\linewidth]{figures/foundation_model/forest_sbs.png}
  \end{subfigure}
  \caption{\textbf{Real (left) and \taeShort reconstructed (right) videos.} The \taeShort compresses the video by a factor of $8\times$ across each of the three spatiotemporal dimensions.
  We observe that the reconstructions from the \taeShort maintain visual detail present in the original videos. %
    }
  \label{fig:tae_examples}
\end{figure}

\noindent\textbf{Quantitative metrics.}
\cref{tab:tae_main_results} compares our \taeShort against a baseline frame-wise autoencoder that does not perform any temporal compression.
Our baseline also produces a 8 channel latent space as is standard for autoencoders used for frame-wise encoding in prior work~\citep{sdvideo,emuvideo2023}.
On video data, we observe that the \taeShort achieves a competitive performance to the frame-wise encoder while achieving a $8\times$ higher temporal compression.
On images, the \taeShort outperforms the frame-wise model, an improvement that can be attributed to the increased channel size of the latent (8 \vs 16)~\citep{dai2023emu}.

\begin{table}[t!]
    \centering
    \begin{tabular}{ccccccc}
        \toprule
         & \multicolumn{3}{c}{Video (512 px)} & \multicolumn{3}{c}{Image (512 px)} \\
         & $\text{SSIM}\uparrow$ & $\text{PSNR}\uparrow$ & $\text{FID}\downarrow$ & $\text{SSIM}\uparrow$ & $\text{PSNR}\uparrow$ & $\text{FID}\downarrow$  \\
         \midrule
         Frame-wise AE & 0.9348 & 34.11 & 0.9352 & 0.8877 & 30.83 & 1.7588 \\
        \taeShort
         & 0.9093 & 32.25 & 1.4872 & 0.9231 & 32.16 & 0.2873 \\
         \bottomrule
    \end{tabular}
    \caption{\textbf{\taeShort reconstruction metrics comparison} between an frame-wise autoencoder and our \taeShort model. We observe that the \taeShort achieves comparable performance to a frame-wise autoencoder for video reconstruction while achieving a $8\times$ higher compression.}
    \label{tab:tae_main_results}
\end{table}



\subsubsection{TAE Ablations}

We now perform a series of ablation experiments for the design choices in training our \taeShort model.

\noindent \textbf{Baseline setting for ablations.}
For simplicity, we use a \taeShort model with a smaller $4\times$ compression ratio that produces a 8-channel latent space.


\noindent\textbf{2.5D \vs 3D attention \& convolutions.}
We compare using 2.5D, \ie, 2D spatial attention/convolutions followed by 1D temporal attention/convolutions to using 3D spatiotemporal attention/convolutions in the \taeShort.
In~\cref{tab:tae_3d_ablation}, we observe that the 3D spatiotemporal attention leads to slightly better reconstruction metrics.
However, we found that this improvement was not large enough to justify the larger memory and compute costs associated with a fully 3D model compared to a 2.5D model.
Thus, we use 2.5D for our \taeShort.


\begin{table}[h!]
    \centering
    \begin{tabular}{ccccccc}
        \toprule
         & \multicolumn{3}{c}{Video (256 px)}\\
         & $\text{SSIM}\uparrow$ & $\text{PSNR}\uparrow$ & $\text{FID}\downarrow$   \\
         \midrule
         \taeShort 2.5D & 0.845 & 28.51 & 3.678\\
         \taeShort 3D & 0.852 & 28.80 & 4.715\\
         \bottomrule
    \end{tabular}
    \caption{\textbf{Comparing \taeShort models with 2.5D \vs 3D convolutions and attentions.} We observe that while the 3D models perform slightly better than the 2.5D models, the gap in performance between the two models is small.
    }
    \label{tab:tae_3d_ablation}
\end{table}

\noindent\textbf{Effect of outlier penalty loss.}
We ablate the effect of adding the outlier penalty loss (OPL) from~\cref{sec:tae}. The addition of this loss removes the artifacts from generated and reconstructed videos as seen in~\cref{fig:tae_black_spot}, and improves reconstruction performance.
We first train a baseline model without OPL for 50K iterations.
We then finetune this model with OPL for 10K iterations and compare it to a baseline finetuned without OPL for 20K iterations.
The results, summarized in~\cref{tab:tae_opl_ablation}, suggest that OPL finetuning improves the reconstruction for both images and videos.
\begin{table}[h!]
    \centering
    \begin{tabular}{ccccccc}
        \toprule
         & \multicolumn{3}{c}{Video (512 px)} & \multicolumn{3}{c}{Image (512 px)} \\
         & $\text{SSIM}\uparrow$ & $\text{PSNR}\uparrow$ & $\text{FID}\downarrow$ & $\text{SSIM}\uparrow$ & $\text{PSNR}\uparrow$ & $\text{FID}\downarrow$  \\
         \midrule
         \taeShort w/o OPL finetune & 0.897 & 31.11 & 1.389 & 0.845 & 28.25 & 0.568 \\
         \taeShort w/ OPL finetune & 0.910 & 31.93 & 1.241 & 0.862 & 29.26 & 0.614 \\
         \bottomrule
    \end{tabular}
    \caption{\textbf{Effect of OPL on \tae reconstructions}. We evaluate the effect of OPL finetuning on the reconstruction quality and observe that OPL improves both image and video metrics.}
    \label{tab:tae_opl_ablation}
\end{table}








